\documentclass[10pt]{extarticle}
\usepackage[a4paper,margin=1.5cm]{geometry}
\usepackage[T1]{fontenc}
\usepackage{amsmath,amssymb}
\usepackage{graphicx}
\usepackage[export]{adjustbox}
\usepackage{multicol}
\usepackage{enumitem}
\usepackage{booktabs,array,tabularx}
\setlength{\columnsep}{0.7cm}
\setlength{\parindent}{0pt}
\setlength{\parskip}{3pt}
\pagestyle{empty}
\makeatletter
\newcommand{\qfigdef}[4]{\@namedef{qf@file@#1}{#2}\@namedef{qf@w@#1}{#3}\@namedef{qf@wide@#1}{#4}}
\newcommand{\qf@check}[1]{\@ifundefined{qf@file@#1}{\PackageError{score2}{Unknown figure id #1}{}}{}}
\newcommand{\qf@img}[4][]{\includegraphics[width=#2\dimexpr\@nameuse{qf@w@#4}pt\relax,max width=#3,max height=0.45\textheight,#1]{\@nameuse{qf@file@#4}}}
\newif\ifqf@full
\newif\ifqf@boxed
\newsavebox{\qf@box}
\newcommand{\qf@open}{\global\qf@fulltrue\par\global\qf@boxedtrue
  \setbox\qf@box\vbox\bgroup\hsize\textwidth\linewidth\textwidth}
\newcommand{\qf@emit}[1]{\par\egroup\global\qf@boxedfalse
  \ifdim\dimexpr\ht\qf@box+\dp\qf@box\relax<#1\textheight\noindent\usebox{\qf@box}\par\else\unvbox\qf@box\fi}
\newcommand{\qfullsplit}{\ifqf@boxed\qf@emit{0.85}\fi}
\newcommand{\qf@close}{\ifqf@boxed\qf@emit{0.45}\fi\global\qf@fullfalse}
\newcommand{\qfullbegin}{\end{multicols}\qf@open}
\newcommand{\qfullend}{\qf@close\begin{multicols}{2}}
\newcommand{\qfullstart}{\qf@open}
\newcommand{\qfulldone}{\qf@close}
\newcommand{\qfig}[2][1]{\qf@check{#2}%
  \ifqf@full\qf@blk{#1}{#2}\else
  \ifnum\@nameuse{qf@wide@#2}=1 \expandafter\@firstoftwo\else\expandafter\@secondoftwo\fi
  {\end{multicols}\begin{center}\qf@img{#1}{\textwidth}{#2}\end{center}\begin{multicols}{2}}%
  {\qf@blk{#1}{#2}}\fi}
\newcommand{\qf@blk}[2]{\par\smallskip{\centering\qf@img{#1}{\linewidth}{#2}\par}\smallskip}
\newcommand{\qfigi}[2][1]{\qf@check{#2}\qf@img[valign=c]{#1}{\linewidth}{#2}}
\makeatother
\newcommand{\papersec}[1]{\par\bigskip{\centering\Large\bfseries PAPER-#1\par}\medskip}
\newcommand{\partsec}[2]{\par\medskip\noindent{\large\bfseries PART-#1: #2}\par\smallskip}
\newcommand{\sect}[1]{\par\smallskip\noindent\textbf{#1}\par}
\newcommand{\qn}[1]{\par\smallskip\noindent\textbf{#1.}~}
\newcommand{\paraq}[2]{\par\smallskip\noindent\textbf{Paragraph for Questions #1 and #2}\par}
\newenvironment{qoptions}{\begin{itemize}[leftmargin=1.8em,itemsep=1pt,topsep=2pt,label={}]}{\end{itemize}}
\newcommand{\opt}[2]{\item[(#1)] #2}
% figure registry (generated by score2_tools.py finalize)
\qfigdef{p1_phy_s1_q01_1}{SRG (Enthusiast) - Test 36 (25 Aug 2024)-img/p1_phy_s1_q01_1.png}{60.5}{0}
\qfigdef{p1_phy_s1_q02_1}{SRG (Enthusiast) - Test 36 (25 Aug 2024)-img/p1_phy_s1_q02_1.png}{130.3}{0}
\qfigdef{p1_phy_s1_q03_1}{SRG (Enthusiast) - Test 36 (25 Aug 2024)-img/p1_phy_s1_q03_1.png}{160.2}{0}
\qfigdef{p1_phy_s2_q05_1}{SRG (Enthusiast) - Test 36 (25 Aug 2024)-img/p1_phy_s2_q05_1.png}{123.7}{0}
\qfigdef{p1_phy_s2_q06_1}{SRG (Enthusiast) - Test 36 (25 Aug 2024)-img/p1_phy_s2_q06_1.png}{102.4}{0}
\qfigdef{p1_phy_s2_q07_1}{SRG (Enthusiast) - Test 36 (25 Aug 2024)-img/p1_phy_s2_q07_1.png}{66.6}{0}
\qfigdef{p1_phy_s3_q03_1}{SRG (Enthusiast) - Test 36 (25 Aug 2024)-img/p1_phy_s3_q03_1.png}{115.6}{0}
\qfigdef{p1_phy_s3_q04_1}{SRG (Enthusiast) - Test 36 (25 Aug 2024)-img/p1_phy_s3_q04_1.png}{107.0}{0}
\qfigdef{p1_phy_s4_q05_1}{SRG (Enthusiast) - Test 36 (25 Aug 2024)-img/p1_phy_s4_q05_1.png}{107.5}{0}
\qfigdef{p1_phy_s4_q06_1}{SRG (Enthusiast) - Test 36 (25 Aug 2024)-img/p1_phy_s4_q06_1.png}{117.5}{0}
\qfigdef{p1_phy_s4_q07_1}{SRG (Enthusiast) - Test 36 (25 Aug 2024)-img/p1_phy_s4_q07_1.png}{131.5}{0}
\qfigdef{p1_phy_s4_q08_1}{SRG (Enthusiast) - Test 36 (25 Aug 2024)-img/p1_phy_s4_q08_1.png}{90.4}{0}
\qfigdef{p1_phy_s4_q09_1}{SRG (Enthusiast) - Test 36 (25 Aug 2024)-img/p1_phy_s4_q09_1.png}{91.8}{0}
\qfigdef{p1_phy_s4_q10_1}{SRG (Enthusiast) - Test 36 (25 Aug 2024)-img/p1_phy_s4_q10_1.png}{77.3}{0}
\qfigdef{p1_chem_s1_q01_1}{SRG (Enthusiast) - Test 36 (25 Aug 2024)-img/p1_chem_s1_q01_1.png}{177.6}{0}
\qfigdef{p1_chem_s1_q01_A}{SRG (Enthusiast) - Test 36 (25 Aug 2024)-img/p1_chem_s1_q01_A.png}{51.5}{0}
\qfigdef{p1_chem_s1_q01_B}{SRG (Enthusiast) - Test 36 (25 Aug 2024)-img/p1_chem_s1_q01_B.png}{50.8}{0}
\qfigdef{p1_chem_s1_q01_C}{SRG (Enthusiast) - Test 36 (25 Aug 2024)-img/p1_chem_s1_q01_C.png}{53.2}{0}
\qfigdef{p1_chem_s1_q01_D}{SRG (Enthusiast) - Test 36 (25 Aug 2024)-img/p1_chem_s1_q01_D.png}{52.4}{0}
\qfigdef{p1_chem_s1_q02_1}{SRG (Enthusiast) - Test 36 (25 Aug 2024)-img/p1_chem_s1_q02_1.png}{145.5}{0}
\qfigdef{p1_chem_s1_q03_1}{SRG (Enthusiast) - Test 36 (25 Aug 2024)-img/p1_chem_s1_q03_1.png}{216.7}{0}
\qfigdef{p1_chem_s2_q05_B}{SRG (Enthusiast) - Test 36 (25 Aug 2024)-img/p1_chem_s2_q05_B.png}{177.1}{0}
\qfigdef{p1_chem_s3_para1_2_1}{SRG (Enthusiast) - Test 36 (25 Aug 2024)-img/p1_chem_s3_para1_2_1.png}{94.2}{0}
\qfigdef{p1_chem_s3_para3_4_1}{SRG (Enthusiast) - Test 36 (25 Aug 2024)-img/p1_chem_s3_para3_4_1.png}{52.4}{0}
\qfigdef{p1_chem_s3_q04_1}{SRG (Enthusiast) - Test 36 (25 Aug 2024)-img/p1_chem_s3_q04_1.png}{96.8}{0}
\qfigdef{p1_chem_s4_q06_1}{SRG (Enthusiast) - Test 36 (25 Aug 2024)-img/p1_chem_s4_q06_1.png}{95.8}{0}
\qfigdef{p1_chem_s4_q07_1}{SRG (Enthusiast) - Test 36 (25 Aug 2024)-img/p1_chem_s4_q07_1.png}{286.3}{1}
\qfigdef{p1_chem_s4_q08_1}{SRG (Enthusiast) - Test 36 (25 Aug 2024)-img/p1_chem_s4_q08_1.png}{332.6}{1}
\qfigdef{p1_chem_s4_q10_1}{SRG (Enthusiast) - Test 36 (25 Aug 2024)-img/p1_chem_s4_q10_1.png}{268.4}{1}

\begin{document}
\begin{center}{\Large\bfseries SRG (Enthusiast) --- Test 36 (25 Aug 2024)}\end{center}
\begin{multicols}{2}
\partsec{1}{PHYSICS}
\sect{Section-I(i)}
\qn{1} A circular conducting loop of radius $R$ carries a current $I$. Another straight infinite conductor carrying current $I$ passes through the diameter of this loop as shown in the figure. The magnitude of force exerted by the straight conductor on the loop is:
\qfig{p1_phy_s1_q01_1}
\begin{qoptions}
\opt{A}{$\pi\mu_0I^2$}
\opt{B}{$\mu_0I^2$}
\opt{C}{$\dfrac{\mu_0I^2}{2\pi}$}
\opt{D}{$\dfrac{\mu_0I^2}{\pi}$}
\end{qoptions}
\qn{2} A long-insulated copper wire is closely wound as a spiral of `$N$' turns. The spiral has inner radius `$a$' and outer radius `$b$'. The spiral lies in the X-Y plane and a steady current flows through the wire. The Z-component of the magnetic field at the center of the spiral is.
\qfig{p1_phy_s1_q02_1}
\begin{qoptions}
\opt{A}{$\dfrac{\mu_0NI}{2(b-a)}\ell n\left(\dfrac ba\right)$}
\opt{B}{$\dfrac{\mu_0NI}{2(b-a)}\ell n\left(\dfrac{b+a}{b-a}\right)$}
\opt{C}{$\dfrac{\mu_0NI}{2b}\ell n\left(\dfrac ba\right)$}
\opt{D}{$\dfrac{\mu_0NI}{2b}\ell n\left(\dfrac{b+a}{b-a}\right)$}
\end{qoptions}
\qn{3} A thin plate separates two liquids of coefficients of viscosity $\eta$ and $4\eta$ kept between two fixed plates as shown. If plate has to be pulled by applying minimum force then $\dfrac{d_2}{d_1}$ is: [Assume steady flow].
\qfig{p1_phy_s1_q03_1}
\begin{qoptions}
\opt{A}{1}
\opt{B}{2}
\opt{C}{3}
\opt{D}{4}
\end{qoptions}
\qn{4} The capillaries of radius $r$ and $2r$ and lengths $\ell$ and $2\ell$ respectively are joined in series. The streamline flow is maintained in capillaries. The pressure difference across first and second capillary is $P_1$ and $P_2$ respectively. Find $P_1/P_2$.
\begin{qoptions}
\opt{A}{$\dfrac14$}
\opt{B}{4}
\opt{C}{$\dfrac18$}
\opt{D}{8}
\end{qoptions}
\sect{Section-I(ii)}
\qn{5} A small solid ball of density $\rho$ is held inside at point A of a cubical container of side $L$, filled with an ideal liquid of density $4\rho$ as shown. The container is closed at the top. Now, if the container starts moving with constant acceleration $a$ horizontally and the ball is released from point A, simultaneously, then:
\qfig{p1_phy_s2_q05_1}
\begin{qoptions}
\opt{A}{For the ball to hit top of container at end Q, $a=3g$}
\opt{B}{For the ball to hit top of container at end Q, $a=2g$}
\opt{C}{Ball hits the top of container at end Q after a time $t=\sqrt{\dfrac{L}{3g}}$}
\opt{D}{Ball hits the top of container at end Q after a time $t=\sqrt{\dfrac{L}{2g}}$}
\end{qoptions}
\qn{6} An electron is at P at $t=0$. It is circulating in anti-clockwise direction with a constant angular speed $\omega$ along the shown circular path. Magnetic field at Q($CQ=2R$, where $R$ is radius of the circle) will be recorded as zero at times.
\qfig{p1_phy_s2_q06_1}
\begin{qoptions}
\opt{A}{$\dfrac{\pi}{3\omega}$}
\opt{B}{$\dfrac{3\pi}{3\omega}$}
\opt{C}{$\dfrac{5\pi}{3\omega}$}
\opt{D}{$\dfrac{7\pi}{3\omega}$}
\end{qoptions}
\qn{7} Consider a cube of side `$a$' as shown. Eight point charges are placed at the corners. The cube is rotated about the axis with constant angular velocity `$\omega$':
\qfig{p1_phy_s2_q07_1}
\begin{qoptions}
\opt{A}{Net magnetic field at the centre of cube is zero}
\opt{B}{Net magnetic field at the centre of cube is $\dfrac{\sqrt2\mu_0q\omega}{\pi a}$}
\opt{C}{Net magnetic field at the centre of cube is $\dfrac{8}{3\sqrt3}\dfrac{\mu_0q\omega}{\pi a}$}
\opt{D}{If polarity of any four charges are reversed, then magnetic field at the centre of cube will be zero.}
\end{qoptions}
\sect{Section-II(i)}
\paraq{1}{2}
A charge particle of mass $m$ and charge $q$ is projected on a rough horizontal XY plane. Both electric and magnetic fields are given by $\vec E=-10\hat k$N/C and $\vec B=-5\hat k$ Tesla are present in the region. The particle enters into the magnetic field at (4, 0, 0) m with a velocity $50\hat J$ m/sec. The particle starts into a curved path on the plane. If coefficient of friction $\mu=\dfrac13$ between particle and plane, then ($qE=2mg$, $g=10\,\mathrm{m/s^2}$).
\qn{1} Radius of curvature of the path followed by particle, initially, is \ldots\ldots\ldots\ldots meters.
\qn{2} The time after which particle comes to rest, is \ldots\ldots\ldots\ldots secs.
\paraq{3}{4}
A wooden cylinder of diameter $4r$, height $H$ and density $\rho/3$ is kept on a hole of diameter $2r$ of a tank, filled with liquid of density $\rho$ as shown in the figure.
\qn{3} If level of the liquid starts decreasing slowly when the level of liquid is at a height $h_1$ above the cylinder the block starts moving up. The value of $h_1$ at which the block will rise is $\dfrac{\alpha H}{3}$. Find the value of $\alpha$.
\qfig{p1_phy_s3_q03_1}
\qn{4} The block in the above question is maintained at the position by external means and the level of liquid is lowered. The height $h_2$ when this external force reduces to zero is $\dfrac{4H}{\beta}$. Find the value of $\beta$.
\qfig{p1_phy_s3_q04_1}
\sect{Section-II(ii)}
\qn{5} A very long straight conductor has a circular cross-section of radius $R$ and carries a current density $J$. Inside the conductor there is a cylindrical hole of radius $a$ whose axis is parallel to the axis of the conductor and a distance $b$ from it. Let the $z$-axis be the axis of the conductor, and let the axis of the hole be at $x=b$. Find the $x$ component of magnetic field on the $y$--axis at $y=2R$. If your answer is $B_x=\mu_0JR\left(\dfrac1A-\dfrac{a^2}{BR^2+b^2}\right)$ fill value of $|A|+|B|$.
\qfig{p1_phy_s4_q05_1}
\qn{6} Find the value of magnetic dipole moment (in $\mathrm{Am^2}$) of the following circuit the current distribution is shown in the diagram.
\qfig{p1_phy_s4_q06_1}
\qn{7} A symmetric star shaped conducting wire loop is carrying a steady state current $I$ as shown in the figure. The distance between the diametrically opposite vertices of the star is $4a$. The magnitude of the magnetic field at the center of the loop is $\dfrac{\alpha\mu_0I\left(\sqrt\alpha-1\right)}{\beta\,\pi\,a}$. Find the value of $3\alpha+4\beta$.
\qfig{p1_phy_s4_q07_1}
\qn{8} A ring of mass $m$ and radius $r$ is rotated in uniform magnetic field $B$ which is perpendicular to the plane of the loop with constant angular velocity $\omega_0$. The tension developed in the ring is $\dfrac{r}{\beta\pi}\left(m\omega_0^{\alpha}+\delta\pi iB\right)$. If there is a current $i$ in the ring. Current and rotation both are clockwise. Find the value of $\sqrt{2\alpha}+3\beta^2+4\delta$. [Axis of rotation is perpendicular to plane and passes through centre].
\qfig{p1_phy_s4_q08_1}
\qn{9} Two immiscible liquids are filled in conical flask as shown in figure. The area of cross section is shown, a small hole of area $a$ is made in lower end of cone. The speed of liquid flow from hole is $\sqrt{\dfrac{3gh}{1-\dfrac{\beta a^2}{32A^2}}}$. Find the value of $\beta$.
\qfig{p1_phy_s4_q09_1}
\qn{10} Water shoots out of a pipe and nozzle as shown in the figure. The cross - sectional area for the tube at point A is four times that of the nozzle outlet. The pressure of water at point A is $41\times10^3\ \mathrm{Nm^{-2}}$ (gauge). If the height `$h$' above the nozzle to which water jet will shoot is $x/10$ meter than $x$ is. (Neglect all the losses occurred in the above process) [ $g=10$ m/s$^2$].
\qfig{p1_phy_s4_q10_1}
\partsec{2}{CHEMISTRY}
\sect{Section-I(i)}
\qn{1}
\qfig{p1_chem_s1_q01_1}
(P) will be :
\begin{qoptions}
\opt{A}{\qfigi{p1_chem_s1_q01_A}}
\opt{B}{\qfigi{p1_chem_s1_q01_B}}
\opt{C}{\qfigi{p1_chem_s1_q01_C}}
\opt{D}{\qfigi{p1_chem_s1_q01_D}}
\end{qoptions}
\qn{2}
\qfig{p1_chem_s1_q02_1}
Select the correct statement.
\begin{qoptions}
\opt{A}{X will give effervescence of $\mathrm{CO_2}$ with $\mathrm{NaHCO_3}$}
\opt{B}{X will produce alkane when treated with Zn / Hg + HCl}
\opt{C}{X will produce symmetrical diketone when heated with sodalime}
\opt{D}{All are correct}
\end{qoptions}
\qn{3} Consider the following compounds
\qfig{p1_chem_s1_q03_1}
Which option satisfy for the above two?
\begin{qoptions}
\opt{A}{Both can show geometrical as well as optical isomerism}
\opt{B}{Both can show optical isomerism}
\opt{C}{Both are chiral but only $2^{\rm nd}$ can show geometrical isomerism}
\opt{D}{Both can show geometrical isomerism}
\end{qoptions}
\qn{4} Buta-1, 3-diene can be converted into But-2-ene by :
\begin{qoptions}
\opt{A}{$\mathrm{LiAlH_4}$/Ether}
\opt{B}{Na/Liquid $\mathrm{NH_3}$}
\opt{C}{Zn-Hg/HCl}
\opt{D}{(i) $\mathrm{BH_3\cdot THF}$ (ii) $\mathrm{AgNO_3}$}
\end{qoptions}
\sect{Section-I(ii)}
\qn{5}
\qfig{p1_chem_s2_q05_B}
Select the correct statement:
\begin{qoptions}
\opt{A}{(D) is $\mathrm{CH_2OH-COO^{(-)}}$}
\opt{B}{(B) \& (E) are same compounds}
\opt{C}{(A) is glyoxal}
\opt{D}{(E) on oxidation gives (B)}
\end{qoptions}
\qn{6} Among the following the false statements are.
\begin{qoptions}
\opt{A}{In $\mathrm{[Co\,Br\,Cl(en)_2]^+}$ geometrical isomerism exist, while optical isomerism does not exist}
\opt{B}{Potassium aquadicyanosuperoxo chromate (III) is the IUPAC name of $\mathrm{K_2[Cr(CN)_2\,O_2(O_2)\,(H_2O)]}$}
\opt{C}{There are 3 geometrical and 15 stereo isomers possible for $\mathrm{[Pt(NO_2)\,(NH_3)\,(NH_2OH)\,(Py)]^+}$ and $\mathrm{[Pt\,Br\,Cl(NO_2)_2\,(NH_3)\,(Py)]}$ respectively}
\opt{D}{Cis and Trans forms are not diastereomers to each other}
\end{qoptions}
\qn{7} Buta 1, 3 diene $+\underset{(1\mathrm{mol})}{\mathrm{HBr}}\to\mathrm{A}+\mathrm{B}$\\
A \& B are primary and secondary allyl bromides. Select the true statement/s:
\begin{qoptions}
\opt{A}{A is thermodynamically controlled product}
\opt{B}{B is thermodynamically controlled product}
\opt{C}{A is kinetically controlled product}
\opt{D}{B is kinetically controlled product}
\end{qoptions}
\sect{Section-II(i)}
\paraq{1}{2}
An organic compound (A) is highly volatile. (A) On reaction with HI gives two products (B) \& (C). (C) gives red colour in victor Mayer test but when (C) is heated with $\mathrm{H^+}$ followed by dilute sulphuric acid, it gives (D) which immediately reacts with Anhydrous $\mathrm{ZnCl_2}$ + HCl. (B) on treatment with $\mathrm{C_2H_5OH}$ + KOH followed by $\mathrm{O_3}$ / $\mathrm{Me_2S}$ gives (E). (E) can gives haloform test. (E) When heated with NaOH gives (F). (F) Also gives haloform test \& can decolourise $\mathrm{Br_2}$ water.
\qfig{p1_chem_s3_para1_2_1}
\qn{1} Degree of unsaturation of F is \underline{\hspace{1cm}}
\qn{2} Molecular weight of D is \underline{\hspace{1cm}}
\paraq{3}{4}
An organic compound (A) containing six membered ring gives (B) when treated with KOH + $\mathrm{C_2H_5OH}$. (B) Gives (C) when subjected to heat with $\mathrm{KMnO_4}$. (C) Releases $\mathrm{CO_2}$ with $\mathrm{NaHCO_3}$ \& gives (D) \& (F) when treated with [CaO / heat] \& [$\mathrm{CH_2N_2}$/heat] respectively. (D) gives (E) when heated with $\mathrm{Ca(OH)_2}$. (E) further gives (D) when treated with [$\mathrm{N_2H_4}$/heat] followed by ozonolysis. On the other hand (F) gives (G) when treated with EtONa/EtOH. (G) Also gives (D) by [$\mathrm{H_3O^+}$/ heat].
\qfig{p1_chem_s3_para3_4_1}
\qn{3} Number of oxygen atoms present in compound G is \underline{\hspace{1cm}}
\qn{4}
\qfig{p1_chem_s3_q04_1}
Number of hyperconjugated hydrogens in one molecule of organic product is \underline{\hspace{1cm}}
\sect{Section-II(ii)}
\qn{5} A, B and C are three complexes of chromium (III) with the empirical formula $\mathrm{H_{12}O_6Cl_3Cr}$. All the three complexes have water and chloride ion as ligands.\\
Complex A does not react with concentrated $\mathrm{H_2SO_4}$, whereas complexes B and C lose 6.75\% and 13.5\% of their original mass, respectively, on treatment with concentrated $\mathrm{H_2SO_4}$. Sum of Ionisable chlorides in A, B \& C are \underline{\hspace{1cm}}
\qn{6} Number of stereoisomers possible for following compounds are
\qfig{p1_chem_s4_q06_1}
\qfullbegin
\qn{7} Number of compounds that give Cannizaro reaction are :
\qfig{p1_chem_s4_q07_1}
\qn{8} In the below reaction, number of moles of $\mathrm{LiAlH_4}$ consumed per mol of the reactant is `$X$'. Then the value of `$4X$' is
\qfig{p1_chem_s4_q08_1}
\qfullend
\qn{9} If all six positions of octahedral geometry are marked by different colours first, then any one bidentate unsymmetrical ligand having no chiral centre can be introduced by \underline{\hspace{1cm}} number of ways. \textbf{(Divide your answer by 6 and fill OMR sheet)}
\qfullbegin
\qn{10} How many of following reaction(s) give(s) Benzoic acid or its salt as product / one of the the product
\qfig{p1_chem_s4_q10_1}
\qfullend
\partsec{3}{MATHEMATICS}
\sect{Section-I(i)}
\qn{1} \par\smallskip\noindent\resizebox{\linewidth}{!}{$\displaystyle\int\dfrac{e^{\tan^{-1}x}}{\left(1+x^2\right)}\left[\left(\sec^{-1}\sqrt{1+x^2}\right)^2+\cos^{-1}\left(\dfrac{1-x^2}{1+x^2}\right)\right]dx,\ (x>0)$}\par\smallskip\noindent, = (where $C$ is constant of integration)
\begin{qoptions}
\opt{A}{$e^{\tan^{-1}x}.\tan^{-1}x+C$}
\opt{B}{$\dfrac{e^{\tan^{-1}x}.\left(\tan^{-1}x\right)^2}{2}+C$}
\opt{C}{$e^{\tan^{-1}x}.\left[\sec^{-1}\left(\sqrt{1+x^2}\right)\right]^2+C$}
\opt{D}{$e^{\tan^{-1}x}.\left[\operatorname{cosec}^{-1}\left(\sqrt{1+x^2}\right)\right]^2+C$}
\end{qoptions}
\qn{2} For any natural number $m$, \par\smallskip\noindent\resizebox{\linewidth}{!}{$\displaystyle\int\left(x^{7m}+x^{2m}+x^m\right)\left(2x^{6m}+7x^m+14\right)^{1/m}dx$}\par\smallskip\noindent (where $x>0$), equals: Where $C$ is constant of integration.
\begin{qoptions}
\opt{A}{$\dfrac{\left(7x^{7m}+2x^{2m}+x^m\right)^{\frac{m+1}m}}{14(m+1)}+C$}
\opt{B}{$\dfrac{\left(2x^{7m}+14x^{2m}+7x^m\right)^{\frac{m+1}m}}{14(m+1)}+C$}
\opt{C}{$\dfrac{\left(2x^{7m}+7x^{2m}+14x^m\right)^{\frac{m+1}m}}{14(m+1)}+C$}
\opt{D}{$\dfrac{\left(2x^{7m}+7x^{2m}+x^m\right)^{\frac{m+1}m}}{14(m+1)}+C$}
\end{qoptions}
\qn{3} The area of the rectangle formed by the perpendiculars from the centre of the standard ellipse $\dfrac{x^2}{a^2}+\dfrac{y^2}{b^2}=1$ to the tangent and normal at its point whose eccentric angle is $\pi/4$, is: (where $a>b$)
\begin{qoptions}
\opt{A}{$\dfrac{\left(a^2-b^2\right)ab}{a^2+b^2}$}
\opt{B}{$\dfrac{\left(a^2-b^2\right)}{\left(a^2+b^2\right)ab}$}
\opt{C}{$\dfrac{\left(a^2+b^2\right)ab}{\left(a^2-b^2\right)}$}
\opt{D}{$\dfrac{a^2+b^2}{\left(a^2-b^2\right)ab}$}
\end{qoptions}
\qn{4} A conic passes through the point (2, 4) and is such that the segment of any of its tangents at any point contained between the co-ordinate axes is bisected at the point of tangency. Then the foci of the conic are:
\begin{qoptions}
\opt{A}{$\left(2\sqrt2,0\right)$ and $\left(-2\sqrt2,0\right)$}
\opt{B}{$\left(2\sqrt2,2\sqrt2\right)$ and $\left(-2\sqrt2,-2\sqrt2\right)$}
\opt{C}{$(4,4)$ and $(-4,-4)$}
\opt{D}{$\left(4\sqrt2,4\sqrt2\right)$ and $\left(-4\sqrt2,-4\sqrt2\right)$}
\end{qoptions}
\sect{Section-I(ii)}
\qn{5} $\displaystyle\int\sqrt{1+\operatorname{cosec}x}\,dx$ equals: (Where $c$ is integration constant and $x\in(0,\pi/2)$)
\begin{qoptions}
\opt{A}{$2\sin^{-1}\sqrt{\sin x}+c$}
\opt{B}{$\sqrt2\cos^{-1}\sqrt{\cos x}+c$}
\opt{C}{$c-2\sin^{-1}(1-2\sin x)$}
\opt{D}{$\cos^{-1}(1-2\sin x)+c$}
\end{qoptions}
\qn{6} Variable circle is described to pass through point (1, 0) and is tangent to the curve $y=\tan\left(\tan^{-1}x\right)$. The locus of the centre of the circle is a parabola whose:
\begin{qoptions}
\opt{A}{Length of the latus rectum is $2\sqrt2$}
\opt{B}{Axis of symmetry has the equation $x+y=1$}
\opt{C}{Vertex has the co-ordinates (3/4, 1/4)}
\opt{D}{Vertex has the co-ordinates $\left(\dfrac12,0\right)$}
\end{qoptions}
\qn{7} A parallelogram circumscribes the ellipse $\dfrac{x^2}9+\dfrac{y^2}4=1$ and two of its opposite angular points lie on straight lines $x^2=\lambda^2,\ \lambda\ne0$, the locus of the other two vertices is:
\begin{qoptions}
\opt{A}{Ellipse (or circle) if $\lambda>3$}
\opt{B}{Hyperbola if $\lambda>3$}
\opt{C}{Circle if $\lambda=\dfrac{\sqrt5}9$}
\opt{D}{Ellipse (or circle) if $\lambda\in(0,3)$}
\end{qoptions}
\sect{Section-II(i)}
\paraq{1}{2}
$\displaystyle\int\dfrac{dx}{\cos^3x-\sin^3x}=A\tan^{-1}(f(x))+B\ \ln\left|\dfrac{\sqrt2+f(x)}{\sqrt2-f(x)}\right|+C$ where $f(x)=\sin x+\cos x$ and $C$ is integration constant.
\qn{1} The value of $3A$ is
\qn{2} The value of $(3B)^2$ is
\paraq{3}{4}
Consider one side AB of a square ABCD (in order) on the line $y=2x-17$, and the other two vertices C, D on the parabola $y=x^2$.
\qn{3} Minimum intercept of the line CD on $y$-axis, is:
\qn{4} Maximum possible area of the square ABCD is
\sect{Section-II(ii)}
\qn{5} If $\displaystyle\int\dfrac{dx}{x^4\left(1+x^3\right)^2}=a\ln\left|\dfrac{1+x^3}{x^3}\right|+\dfrac b{x^3}+\dfrac c{1+x^3}+d$, then $a+b+c=$\\
(where $d$ is arbitrary constant)
\qn{6} $\displaystyle\int\dfrac{\cos9x+\cos6x}{2\cos5x-1}dx=A\sin4x+B\sin x+C$, then $[A+B]$ is equal to:\\
(Where $C$ is constant of integration; $[\cdot]$ is greatest integer function)
\qn{7} The value of $\displaystyle\int\dfrac{\tan x}{\tan^2x+\tan x+1}dx=x-\dfrac2{\sqrt A}\tan^{-1}\left(\dfrac{2\tan x+1}{\sqrt A}\right)+C$ (where $C$ is integration constant). Then the value of $A$ is :
\qn{8} The locus of the point of intersection of two tangents to the parabola $y^2=4ax$ which with the tangent at the vertex form a triangle of constant area $c^2$, is the curve $x^2\left(y^2-4ax\right)=\lambda c^4$, then $\lambda=$
\qn{9} A tangent to the ellipse $x^2+4y^2=4$ meets the ellipse $x^2+2y^2=6$ at P and Q. Let angle between tangents at P and Q of the ellipse $x^2+2y^2=6$ is $\dfrac\pi k$, then $k=$
\qn{10} The chord AC of the parabola $y^2=4ax$ subtends an angle of $90^\circ$ at points B and D on the parabola. If points A, B, C and D are represented by $\left(at_i^2,\ 2at_i\right)$, $i=1,2,3,4$ respectively, then find the value of $\left|\dfrac{t_2+t_4}{t_1+t_3}\right|$.
\end{multicols}
\end{document}
